% ─────────────────────────────────────────────────────────────────────────────
%
% mejora_performance.tex
%
% Nota de ingeniería sobre la optimización de performance del pipeline
% de percepción en Raspberry Pi con RealSense D435i.
%
% ─────────────────────────────────────────────────────────────────────────────

\section{Mejora de performance del pipeline de percepción}
\label{sec:mejora-performance}

Esta nota documenta el proceso de diagnóstico y optimización del pipeline de
percepción ejecutado en Raspberry Pi con una cámara Intel RealSense D435i. El
problema inicial se manifestó como bloqueos intermitentes o sensación de
``cuelgue'' al ejecutar el sistema con la configuración:

\begin{verbatim}
ros2 launch nav_bringup nav.launch.py \
  use_realsense:=true \
  use_hw:=false \
  pipeline_mode:=filtered \
  use_perception:=true \
  use_visual_odometry:=false \
  use_local_mapper:=false \
  use_rviz:=false \
  use_viz:=false
\end{verbatim}

En esta configuración se utiliza la RealSense como fuente de profundidad e
IMU, se ejecuta el pipeline filtrado de obstáculos y se desactiva la odometría
visual. El objetivo fue identificar si la causa de los bloqueos provenía del
driver USB/cámara, del procesamiento de percepción, o de una combinación de
ambos.

\subsection{Configuración inicial y síntomas}
\label{subsec:perf-sintomas}

La cámara se configuró explícitamente en el launch para publicar profundidad a
$640 \times 480$ píxeles y $6\,\mathrm{Hz}$:

\begin{verbatim}
depth_module.depth_profile = 640x480x6
\end{verbatim}

La tasa de $6\,\mathrm{Hz}$ implica un presupuesto temporal de:

\begin{equation}
  T_{\mathrm{frame}} = \frac{1000}{6} \approx 166{,}7\,\mathrm{ms}
\end{equation}

Por lo tanto, cualquier etapa de procesamiento que supere de manera sostenida
ese tiempo comienza a acumular cola de mensajes. El síntoma percibido por el
usuario no necesariamente aparece como una caída del nodo, sino como latencia
creciente: el sistema sigue procesando, pero lo hace sobre frames viejos.

En paralelo aparecieron advertencias del driver de RealSense:

\begin{verbatim}
control_transfer returned error, index: 768,
error: Resource temporarily unavailable, number: 11
\end{verbatim}

y, en algunas corridas, también:

\begin{verbatim}
control_transfer returned error, index: 768,
error: Success, number: 0
\end{verbatim}

Estas advertencias provienen de \texttt{librealsense}/USB y no de los nodos de
percepción. Se observó que se intensificaban al mover la cámara, lo que sugiere
un componente físico asociado al cable, conector, alimentación o ruido en el
bus USB. Sin embargo, las mediciones de performance mostraron que también
existía un cuello de botella algorítmico claro en el pipeline filtrado.

\subsection{Instrumentación utilizada}
\label{subsec:perf-instrumentacion}

Se utilizaron tres fuentes principales de medición:

\begin{enumerate}
  \item Logs de performance de \texttt{depth\_obstacle\_filter}.
  \item Logs de performance de \texttt{obstacle\_grid\_encoder}.
  \item \texttt{ros2 topic delay} sobre el tópico crudo de profundidad:
        \texttt{/camera/camera/depth/image\_rect\_raw}.
\end{enumerate}

El log de \texttt{depth\_obstacle\_filter} separa el tiempo total por frame en
etapas:

\begin{itemize}
  \item \texttt{proj}: retroproyección de la imagen de profundidad a puntos 3D.
  \item \texttt{range}: filtrado de rango y descarte de valores inválidos.
  \item \texttt{ground}: rotación a frame alineado a gravedad, estimación de
        suelo y clasificación suelo/obstáculo/techo.
  \item \texttt{publish}: armado y publicación de mensajes.
  \item \texttt{total}: tiempo total medido por el nodo para procesar el frame.
  \item \texttt{age}: edad del mensaje al finalizar el procesamiento.
\end{itemize}

La edad del mensaje se calcula como:

\begin{equation}
  \mathrm{age} = t_{\mathrm{now}} - t_{\mathrm{stamp}}
\end{equation}

donde $t_{\mathrm{stamp}}$ es el timestamp del mensaje de profundidad. Por lo
tanto, \texttt{age} incluye tanto el retardo de cámara/ROS como el tiempo de
procesamiento del nodo.

Para separar ambos efectos se midió el tópico crudo:

\begin{verbatim}
ros2 topic delay /camera/camera/depth/image_rect_raw
\end{verbatim}

El resultado típico fue:

\begin{verbatim}
average delay: 0.010s - 0.014s
max: 0.071s
\end{verbatim}

Esto indica que la cámara entregaba la imagen con un retardo promedio de apenas
$10$ a $14\,\mathrm{ms}$. En consecuencia, cuando
\texttt{depth\_obstacle\_filter} reportaba \texttt{age} cercano a
$150\,\mathrm{ms}$, la mayor parte correspondía al procesamiento del propio
nodo y no a una cola previa en el driver.

\subsection{Medición inicial del pipeline filtrado}
\label{subsec:perf-medicion-inicial}

La primera medición mostró que \texttt{obstacle\_grid\_encoder} no era el
cuello de botella. Sus tiempos típicos se mantuvieron en torno a:

\begin{table}[H]
\centering
\begin{tabular}{lll}
\toprule
\textbf{Nodo} & \textbf{Métrica} & \textbf{Valor observado} \\
\midrule
\texttt{obstacle\_grid\_encoder} & \texttt{avg\_total} & $9$--$10\,\mathrm{ms}$ \\
\texttt{obstacle\_grid\_encoder} & \texttt{max\_total} & $10$--$13\,\mathrm{ms}$ \\
\texttt{obstacle\_grid\_encoder} & \texttt{avg\_age} & $160$--$185\,\mathrm{ms}$ \\
\texttt{obstacle\_grid\_encoder} & puntos de entrada & $160\,000$--$200\,000$ \\
\bottomrule
\end{tabular}
\caption{Mediciones iniciales del codificador de grilla de obstáculos.}
\label{tab:perf-obstacle-grid-inicial}
\end{table}

En cambio, \texttt{depth\_obstacle\_filter} operaba muy cerca del presupuesto
de $166{,}7\,\mathrm{ms}$ por frame:

\begin{table}[H]
\centering
\begin{tabular}{lll}
\toprule
\textbf{Etapa} & \textbf{Valor típico} & \textbf{Comentario} \\
\midrule
\texttt{proj} & $18$--$29\,\mathrm{ms}$ & Proyección de $307\,200$ píxeles \\
\texttt{range} & $9$--$14\,\mathrm{ms}$ & Filtrado de profundidad \\
\texttt{ground} & $109$--$119\,\mathrm{ms}$ & Etapa dominante \\
\texttt{total} & $142$--$152\,\mathrm{ms}$ & Muy cerca del presupuesto \\
\texttt{max\_total} & $162$--$181\,\mathrm{ms}$ & Supera el período de frame \\
\texttt{age} & $150$--$180\,\mathrm{ms}$ & Latencia final del frame \\
\bottomrule
\end{tabular}
\caption{Medición inicial de \texttt{depth\_obstacle\_filter} con resolución completa.}
\label{tab:perf-depth-filter-inicial}
\end{table}

En esa configuración se procesaban:

\begin{verbatim}
pts proj = 307200
valid    = 260k - 274k
obs      = 160k - 180k
\end{verbatim}

El desglose interno de \texttt{GroundEstimator} mostraba:

\begin{table}[H]
\centering
\begin{tabular}{lll}
\toprule
\textbf{Subetapa} & \textbf{Valor observado} & \textbf{Interpretación} \\
\midrule
\texttt{GE total} & $74$--$86\,\mathrm{ms}$ & Costo total del estimador \\
\texttt{voxel} & $53$--$61\,\mathrm{ms}$ & Costo dominante inicial \\
\texttt{ransac} & $\sim 1\,\mathrm{ms}$ & Bajo debido al downsampling \\
\texttt{refine} & $12$--$16\,\mathrm{ms}$ & Refinamiento sobre inliers \\
\bottomrule
\end{tabular}
\caption{Costo interno inicial del estimador de suelo.}
\label{tab:perf-ground-estimator-inicial}
\end{table}

La conclusión inicial fue que el voxel grid, introducido para reducir la nube
antes de RANSAC, se había convertido en un costo significativo. Esto ocurrió
porque el RANSAC ya no se ejecutaba sobre toda la nube sino sobre una banda
geométrica cercana al máximo de $Y$ medido, por lo que el beneficio del voxel
global dejó de compensar su costo.

\subsection{Desactivación configurable del voxel grid}
\label{subsec:perf-voxel-disable}

Se agregó el parámetro:

\begin{verbatim}
enable_voxel_filter: false
\end{verbatim}

Cuando este parámetro está deshabilitado, \texttt{GroundEstimator} saltea el
voxel grid y usa la nube original para la búsqueda de la banda de suelo. La
medición mostró una mejora en promedio:

\begin{table}[H]
\centering
\begin{tabular}{llll}
\toprule
\textbf{Métrica} & \textbf{Con voxel} & \textbf{Sin voxel} & \textbf{Resultado} \\
\midrule
\texttt{voxel} & $53$--$61\,\mathrm{ms}$ & $\sim 0{,}7\,\mathrm{ms}$ & Mejora grande \\
\texttt{GE total} & $80$--$86\,\mathrm{ms}$ & $53$--$79\,\mathrm{ms}$ & Mejora promedio \\
\texttt{total} & $148$--$152\,\mathrm{ms}$ & $116$--$145\,\mathrm{ms}$ & Mejora promedio \\
\texttt{ransac} & $\sim 1\,\mathrm{ms}$ & $7$--$18\,\mathrm{ms}$ & Aumenta \\
\bottomrule
\end{tabular}
\caption{Efecto de desactivar el voxel grid global.}
\label{tab:perf-voxel-disable}
\end{table}

Sin embargo, también aparecieron picos severos:

\begin{verbatim}
ransac   = 296ms
GE total = 357ms
total    = 454ms
age      = 1326ms
\end{verbatim}

El motivo fue que, al quitar el voxel grid, el conteo de inliers de RANSAC podía
quedar trabajando sobre bandas con decenas o cientos de miles de puntos. Esto
eliminó el costo fijo del voxel, pero dejó sin acotar el peor caso de RANSAC.

\subsection{Limitación de puntos de banda en RANSAC}
\label{subsec:perf-ransac-limit}

Para acotar el peor caso se agregó el parámetro:

\begin{verbatim}
ransac_max_band_points: 5000
\end{verbatim}

El algoritmo sigue construyendo la banda candidata de suelo de la misma forma,
pero si la banda supera ese máximo, se submuestrea de manera determinística
antes de contar inliers. Esto mantiene estable la latencia del conteo sin
reintroducir el voxel grid global.

Con el limitador activo y \texttt{enable\_voxel\_filter: false}, las mediciones
se estabilizaron:

\begin{table}[H]
\centering
\begin{tabular}{lll}
\toprule
\textbf{Métrica} & \textbf{Valor típico con limitador} & \textbf{Comentario} \\
\midrule
\texttt{ransac} & $13$--$25\,\mathrm{ms}$ & Sin picos de cientos de ms \\
\texttt{GE total} & $72$--$86\,\mathrm{ms}$ & Estable \\
\texttt{total} & $136$--$151\,\mathrm{ms}$ & Bajo el presupuesto en promedio \\
\texttt{max\_total} & $152$--$188\,\mathrm{ms}$ & Algunos picos cerca del límite \\
\texttt{age} & $154$--$176\,\mathrm{ms}$ & Sin cola acumulada severa \\
\bottomrule
\end{tabular}
\caption{Efecto del limitador de puntos de banda en RANSAC.}
\label{tab:perf-ransac-limit}
\end{table}

Esta etapa confirmó que la combinación ``sin voxel global + RANSAC acotado''
era mejor que el pipeline original con voxel, pero el sistema aún trabajaba
relativamente cerca del límite de $6\,\mathrm{Hz}$.

\subsection{Submuestreo temprano de profundidad}
\label{subsec:perf-depth-stride}

Como no fue posible utilizar resoluciones menores de la RealSense de manera
confiable en el entorno probado, se mantuvo la cámara en $640 \times 480$ a
$6\,\mathrm{Hz}$ y se introdujo submuestreo de píxeles dentro del nodo:

\begin{verbatim}
depth_pixel_stride: 2
\end{verbatim}

El parámetro se aplicó en \texttt{DepthProjector}, es decir, antes de construir
la nube de puntos. Esto reduce el costo de todas las etapas posteriores:

\begin{itemize}
  \item proyección,
  \item filtrado de rango,
  \item rotación a frame alineado,
  \item estimación de suelo,
  \item publicación de nube,
  \item codificación de la grilla.
\end{itemize}

Con \texttt{depth\_pixel\_stride: 2}, la cantidad de muestras baja de:

\begin{equation}
  640 \times 480 = 307\,200
\end{equation}

a aproximadamente:

\begin{equation}
  320 \times 240 = 76\,800
\end{equation}

La mejora observada fue sustancial:

\begin{table}[H]
\centering
\begin{tabular}{llll}
\toprule
\textbf{Métrica} & \textbf{Antes} & \textbf{Con stride 2} & \textbf{Mejora} \\
\midrule
\texttt{pts proj} & $307\,200$ & $76\,800$ & $4\times$ menos puntos \\
\texttt{proj} & $18$--$29\,\mathrm{ms}$ & $5$--$7\,\mathrm{ms}$ & $\sim 3\times$ \\
\texttt{range} & $9$--$14\,\mathrm{ms}$ & $2{,}4$--$3{,}5\,\mathrm{ms}$ & $\sim 3\times$ \\
\texttt{ground} & $100$--$118\,\mathrm{ms}$ & $26$--$35\,\mathrm{ms}$ & $\sim 3$--$4\times$ \\
\texttt{GE total} & $72$--$86\,\mathrm{ms}$ & $19$--$23\,\mathrm{ms}$ & $\sim 4\times$ \\
\texttt{ransac} & $13$--$25\,\mathrm{ms}$ & $1{,}6$--$6{,}2\,\mathrm{ms}$ & Estable \\
\texttt{total} & $136$--$151\,\mathrm{ms}$ & $35$--$50\,\mathrm{ms}$ & $\sim 3\times$ \\
\texttt{age} & $150$--$176\,\mathrm{ms}$ & $52$--$63\,\mathrm{ms}$ & $\sim 3\times$ \\
\bottomrule
\end{tabular}
\caption{Comparación antes y después de aplicar \texttt{depth\_pixel\_stride: 2}.}
\label{tab:perf-depth-stride}
\end{table}

Ejemplos representativos de la medición final:

\begin{verbatim}
pts proj=76800 valid=65276
proj=6.73ms range=3.21ms ground=33.33ms
total=44.77ms age=57.0ms
GE total=22.58ms voxel=0.17ms ransac=3.89ms refine=3.84ms
\end{verbatim}

\begin{verbatim}
pts proj=76800 valid=66769
proj=6.82ms range=3.25ms ground=33.05ms
total=45.26ms age=62.9ms
GE total=21.84ms voxel=0.17ms ransac=2.66ms refine=4.15ms
\end{verbatim}

\begin{verbatim}
pts proj=76800 valid=66246
proj=4.87ms range=2.43ms ground=26.13ms
total=35.09ms age=51.7ms
GE total=19.31ms voxel=0.16ms ransac=6.24ms refine=2.50ms
\end{verbatim}

El codificador de grilla también redujo su costo:

\begin{table}[H]
\centering
\begin{tabular}{lll}
\toprule
\textbf{Métrica} & \textbf{Antes de stride} & \textbf{Con stride 2} \\
\midrule
\texttt{avg\_total} & $8$--$10\,\mathrm{ms}$ & $2{,}4$--$4{,}2\,\mathrm{ms}$ \\
\texttt{max\_total} & $9$--$13\,\mathrm{ms}$ & $2{,}8$--$6{,}7\,\mathrm{ms}$ \\
\texttt{avg\_points\_in} & $140\,000$--$170\,000$ & $28\,000$--$45\,000$ \\
\texttt{avg\_age} & $160$--$180\,\mathrm{ms}$ & $59$--$68\,\mathrm{ms}$ \\
\bottomrule
\end{tabular}
\caption{Efecto del stride sobre \texttt{obstacle\_grid\_encoder}.}
\label{tab:perf-grid-stride}
\end{table}

\subsection{Publicación asíncrona de \texttt{ObstacleCloud}}
\label{subsec:perf-obstaclecloud-async}

Luego de reducir el costo algorítmico de proyección, filtrado y estimación de
suelo, apareció un segundo cuello de botella que no estaba representado de
forma explícita en las primeras mediciones. El log original sólo informaba una
etapa genérica de publicación, pero el tiempo total del callback seguía siendo
mucho mayor que la suma de las etapas principales.

Para aislar este costo se agregó una métrica específica:

\begin{verbatim}
obs_pub = tiempo asociado a preparar/publicar /depth_obstacle_filter/obstacle_cloud
\end{verbatim}

La medición mostró que la serialización y publicación de
\texttt{/depth\_obstacle\_filter/obstacle\_cloud}, implementada como
\texttt{sensor\_msgs/PointCloud2}, era el bloqueo dominante del callback de
profundidad. Un ejemplo representativo antes del cambio fue:

\begin{verbatim}
obs_pub=108.88/202.65ms
total=116.70/211.58ms
\end{verbatim}

y en otras ventanas:

\begin{verbatim}
obs_pub=148.56/348.20ms
total=157.31/356.63ms

obs_pub=157.99/330.19ms
total=166.61/339.07ms
\end{verbatim}

En esas mismas mediciones, el costo real de cómputo del filtro era bajo:

\begin{verbatim}
proj   ~= 2ms
range  ~= 1ms
ground ~= 5ms - 10ms
\end{verbatim}

Por lo tanto, el problema ya no estaba en RANSAC ni en la clasificación de
suelo, sino en que el callback de profundidad quedaba bloqueado mientras se
construía, serializaba y publicaba una nube de puntos. Esto generaba una cola
creciente de frames y aumentaba la edad de los mensajes, aun cuando el cálculo
geométrico ya entraba dentro del presupuesto temporal.

La solución aplicada fue separar la publicación de \texttt{ObstacleCloud} en un
thread dedicado. El callback de profundidad ahora:

\begin{enumerate}
  \item calcula la nube filtrada,
  \item recorta la nube al volumen realmente consumido por
        \texttt{obstacle\_grid\_encoder},
  \item copia el último trabajo de publicación en una cola de tamaño uno,
  \item retorna sin esperar a que finalice la serialización/publicación del
        \texttt{PointCloud2}.
\end{enumerate}

El thread publicador consume esa cola y publica siempre el último frame
disponible. Si el publicador se atrasa, el trabajo pendiente se reemplaza por
el más reciente. Esta política evita acumular frames viejos y prioriza latencia
baja, que es más importante para realimentación háptica que publicar todos los
frames históricos.

La configuración quedó controlada por:

\begin{verbatim}
async_obstacle_cloud_publish: true
prune_obstacle_cloud_for_grid: true
\end{verbatim}

El recorte previo usa los mismos límites espaciales que consume la grilla:

\begin{verbatim}
obstacle_cloud_x_min_m: -1.0
obstacle_cloud_x_max_m:  1.0
obstacle_cloud_y_min_m: -0.20
obstacle_cloud_y_default_max_m: 2.20
obstacle_cloud_y_margin_m: 0.10
\end{verbatim}

Luego del cambio, la métrica \texttt{obs\_pub} pasó a medir principalmente el
costo de preparar y encolar el trabajo, no la publicación bloqueante completa.
Los valores observados bajaron de cientos de milisegundos a pocos
milisegundos:

\begin{table}[H]
\centering
\begin{tabular}{llll}
\toprule
\textbf{Métrica} & \textbf{Antes} & \textbf{Después} & \textbf{Impacto} \\
\midrule
\texttt{obs\_pub promedio} & $108$--$158\,\mathrm{ms}$ & $0{,}5$--$1{,}8\,\mathrm{ms}$ & Elimina el bloqueo \\
\texttt{obs\_pub máximo} & $202$--$348\,\mathrm{ms}$ & $1{,}3$--$2{,}8\,\mathrm{ms}$ & Sin picos críticos \\
\texttt{total promedio} & $116$--$166\,\mathrm{ms}$ & $5$--$16\,\mathrm{ms}$ & Callback liviano \\
\texttt{total máximo} & $211$--$356\,\mathrm{ms}$ & $10$--$29\,\mathrm{ms}$ & Bajo presupuesto \\
\bottomrule
\end{tabular}
\caption{Impacto de publicar \texttt{ObstacleCloud} en un thread separado.}
\label{tab:perf-obstaclecloud-async}
\end{table}

Un log posterior al cambio muestra el nuevo orden de magnitud:

\begin{verbatim}
frames=65 age=277.9ms
proj=1.59/3.71ms range=0.57/1.25ms ground=2.87/8.37ms
obs_pub=0.56/1.32ms total=5.60/10.68ms
\end{verbatim}

\begin{verbatim}
frames=44 age=1472.8ms
proj=2.11/2.59ms range=0.89/1.14ms ground=6.97/13.15ms
obs_pub=1.08/2.06ms total=11.06/17.63ms
\end{verbatim}

El valor de \texttt{age} puede seguir creciendo si la RealSense deja de
entregar frames, por ejemplo ante eventos:

\begin{verbatim}
uvc streamer watchdog triggered on endpoint: 130
\end{verbatim}

pero esos casos quedan claramente separados del costo de procesamiento del
pipeline. Después de esta optimización, cuando aparecen pausas largas la causa
observada está en cámara/USB/driver, no en el callback de percepción.

\subsection{Configuración de RealSense según odometría visual}
\label{subsec:perf-realsense-config}

Durante el diagnóstico se identificó que, cuando
\texttt{use\_visual\_odometry:=false}, no era necesario publicar color ni
alinear profundidad a color. Se modificó el launch principal para que el nodo
RealSense se configure dinámicamente según \texttt{use\_visual\_odometry}.

Con odometría visual desactivada:

\begin{verbatim}
enable_color = false
align_depth.enable = false
depth_module.depth_profile = realsense_depth_profile
\end{verbatim}

Con odometría visual activada:

\begin{verbatim}
enable_color = true
align_depth.enable = true
depth_module.depth_profile = realsense_depth_profile
rgb_camera.color_profile = realsense_color_profile
\end{verbatim}

Esto reduce tráfico USB y carga de CPU en el modo usado para estas mediciones.
El cambio no altera el comportamiento cuando se requiere RGB-D para odometría
visual. Durante las pruebas finales se usaron perfiles de menor resolución,
por ejemplo:

\begin{verbatim}
realsense_depth_profile:=424x240x15
\end{verbatim}

y, para priorizar estabilidad del enlace USB por sobre tasa máxima:

\begin{verbatim}
realsense_depth_profile:=424x240x6
\end{verbatim}

\subsection{Modo alternativo de orientación por IMU legacy}
\label{subsec:perf-imu-legacy}

Se agregó además un selector de fuente de orientación para
\texttt{depth\_obstacle\_filter}:

\begin{verbatim}
orientation_source: nav_odom
\end{verbatim}

o bien:

\begin{verbatim}
orientation_source: imu_legacy
\end{verbatim}

El modo \texttt{nav\_odom} mantiene el comportamiento actual: se lanza
\texttt{nav\_odometry} y \texttt{depth\_obstacle\_filter} toma roll/pitch desde
el tópico \texttt{nav\_odom}. El modo \texttt{imu\_legacy} no lanza
\texttt{nav\_odometry}; en su lugar, \texttt{depth\_obstacle\_filter} se
suscribe al acelerómetro de la RealSense y utiliza \texttt{ImuFilter} junto con
\texttt{GravityAligner}. Este cambio quedó preparado para comparar la
estabilidad de la estimación de altura cuando la orientación proviene sólo del
acelerómetro filtrado.

La motivación para este selector surgió de los rechazos de altura:

\begin{verbatim}
[camera_height] conservada: piso poco confiable
valid=1 quality=1.000 inliers=... tilt=8.7deg
\end{verbatim}

El filtro de altura estaba configurado con:

\begin{verbatim}
height_max_ground_tilt_deg: 8.0
\end{verbatim}

Por lo tanto, planos detectados con inclinación apenas superior a $8^\circ$
eran rechazados. Esto no era un problema de performance, sino de criterio
geométrico y/o de calidad de la orientación usada para alinear la nube.

\subsection{Estado final recomendado}
\label{subsec:perf-estado-final}

La configuración de performance que produjo el mejor comportamiento medido fue:

\begin{verbatim}
depth_pixel_stride: 2
enable_voxel_filter: false
ransac_max_band_points: 5000
prune_obstacle_cloud_for_grid: true
async_obstacle_cloud_publish: true
perf_log_enabled: true
\end{verbatim}

En las pruebas finales se utilizó además un perfil de profundidad reducido:

\begin{verbatim}
realsense_depth_profile:=424x240x15
\end{verbatim}

Con esta configuración, el callback principal de percepción opera con amplio
margen respecto del presupuesto de $66{,}7\,\mathrm{ms}$ correspondiente a
$15\,\mathrm{Hz}$:

\begin{table}[H]
\centering
\begin{tabular}{lll}
\toprule
\textbf{Métrica final} & \textbf{Valor típico} & \textbf{Evaluación} \\
\midrule
\texttt{total} & $5$--$16\,\mathrm{ms}$ & Muy por debajo del límite \\
\texttt{obs\_pub} & $0{,}5$--$1{,}8\,\mathrm{ms}$ & No bloqueante \\
\texttt{GE total} & $1{,}4$--$10{,}9\,\mathrm{ms}$ & Estable, salvo escenas difíciles \\
\texttt{obstacle\_grid} & $0{,}5$--$1{,}5\,\mathrm{ms}$ & No limitante \\
\texttt{depth/grid rate} & $8$--$14\,\mathrm{Hz}$ & Limitado por RealSense/USB \\
\bottomrule
\end{tabular}
\caption{Estado final del pipeline luego de las optimizaciones.}
\label{tab:perf-final}
\end{table}

El problema de performance quedó entonces separado de los problemas restantes:

\begin{enumerate}
  \item Las advertencias USB/UVC de RealSense siguen indicando una posible
        fragilidad física o de alimentación, especialmente ante mensajes como
        \texttt{uvc streamer watchdog triggered on endpoint: 130}.
  \item La estabilidad de la altura depende de la calidad de la orientación y
        del umbral \texttt{height\_max\_ground\_tilt\_deg}. Una prueba futura
        razonable es comparar \texttt{orientation\_source:=nav\_odom} contra
        \texttt{orientation\_source:=imu\_legacy}.
  \item El submuestreo espacial puede perder obstáculos pequeños. Aunque las
        mediciones de performance son claramente mejores, la configuración
        \texttt{depth\_pixel\_stride: 2} debe validarse en escenarios reales de
        navegación y seguridad.
\end{enumerate}

\subsection{Conclusiones}
\label{subsec:perf-conclusiones}

El análisis mostró que el bloqueo inicial no se explicaba únicamente por las
advertencias USB de RealSense. El pipeline filtrado estaba operando al límite
de la Raspberry Pi: con resolución completa, \texttt{depth\_obstacle\_filter}
tardaba típicamente $135$--$150\,\mathrm{ms}$ por frame y ocasionalmente
superaba el período de $166{,}7\,\mathrm{ms}$.

La primera optimización, desactivar el voxel grid global, redujo el promedio
pero introdujo picos de RANSAC. La segunda, limitar la cantidad de puntos de la
banda de RANSAC, estabilizó esos picos. Luego, el submuestreo temprano de la
imagen de profundidad con \texttt{depth\_pixel\_stride: 2} redujo la nube de
entrada de $307\,200$ a $76\,800$ puntos y bajó el costo geométrico del filtro.

Sin embargo, la mejora de mayor impacto apareció al instrumentar por separado
la publicación de \texttt{ObstacleCloud}. Esa medición mostró que
\texttt{obs\_pub} podía consumir $100$--$160\,\mathrm{ms}$ en promedio y
superar los $300\,\mathrm{ms}$ en picos, aun cuando el procesamiento geométrico
ya era barato. Al mover la publicación del \texttt{PointCloud2} a un thread
separado y usar una cola de tamaño uno con política de último frame, el callback
principal bajó a valores típicos de $5$--$16\,\mathrm{ms}$.

Con estos cambios, el sistema dejó de estar limitado por CPU en el pipeline de
percepción. Los problemas restantes pertenecen a dos categorías distintas:
robustez física del enlace RealSense/USB, evidenciada por eventos UVC watchdog,
y estabilidad geométrica de la estimación de suelo/altura.
